\bookchapter{第三十三章：未来展望——手机的下一程}

\chapterintro{手机不会一夜之间消失，但它的“独霸地位”可能被新一代设备分流。这一章我们聊三件事：下一代设备长什么样、手机如何“隐身”为新身份基石、以及那些我们必须提前思考的伦理议题。}

\booksection{一、AR/VR与空间计算：手机之外的“第二桌面”}

下一代设备的竞争已经打响，最有可能分流手机时间的，是AR/VR与空间计算设备：

\begin{itemize}
  \item \textbf{Vision Pro后续}：苹果Vision Pro的发布标志着"空间计算"概念的正式落地，但第一代产品价格（29999元起）和重量仍是门槛。预计2026-2027年会推出更轻、更便宜的平价版Vision，让空间计算从"开发者玩具"走向"早期大众"。届时iPhone可能成为Vision眼镜的算力伴侣与数据中枢。
  \item \textbf{AR眼镜}：给你的眼前开个"第二桌面"，数字信息变得像空气一样无处不在。Meta与Ray-Ban合作的智能眼镜已卖出数百万副，国内Rokid、XREAL、INMO等厂商也在推出消费级AR眼镜。2025-2026年的AR眼镜已能做实时翻译、导航HUD、AI助手对话，但仍受限于光波导显示和续航，距离"取代手机"还有3-5年。
  \item \textbf{智能戒指与可穿戴}：Samsung Galaxy Ring、Oura Ring等智能戒指开始流行，作为手机的"贴身传感器"补充健康监测。手机的角色正从"万能终端"转向"中枢+多个轻量化配件"的组合。
\end{itemize}

\booksection{二、脑机接口：思想成为交互媒介}

思想本身将成为交互的媒介，但这条路比想象中漫长：

\begin{itemize}
  \item \textbf{Neuralink最新进展}：马斯克的Neuralink已于2024年完成首例人体植入，首位受试者Noland Arbaugh通过意念控制电脑鼠标下国际象棋。目前植入式BCI（脑机接口）主要面向瘫痪患者辅助场景，目标是帮助他们重新"操作"数字世界。
  \item \textbf{非侵入式BCI}：多家公司开发非侵入式头环（如Kernel、OpenBCI），无需手术即可读取脑电信号，但精度远低于侵入式。距离"用意念打字"还有相当距离。
  \item \textbf{消费级展望}：脑机接口距离消费级手机交互还有10-20年，但短期可能先用于无障碍辅助、专注力监测等场景。伦理与隐私问题（"我的脑电波算谁的隐私"）将先于技术成熟成为争论焦点。
\end{itemize}

\booksection{三、手机的“隐身”与新身份基石}

手机的另一个未来进化方向或许是“隐身”——化身为无处不在的“幕后推手”，成为你数字身份的唯一基石：

\begin{itemize}
  \item \textbf{云端双因素认证}：瞬间完成"你是你"的认证，登录任何服务无需手动输验证码。
  \item \textbf{数字遗产成为出厂设置}：新机激活时预设数字资产继承方案，让"身后事"成为出厂默认。
  \item \textbf{从"设备中心"到"身份中心"}：你的云端个人数据保险柜才是中心，硬件只是访问入口。换机、换设备，身份与数据始终跟着你。
\end{itemize}

\booksection{四、伦理思考：我们，还是手机的主人吗？}

一个终极拷问始终存在：我们是被技术赋能的“神”，还是被算法喂养的“缸中之脑”？

\begin{itemize}
  \item \textbf{AI伦理}：当端侧AI比你自己更懂你的偏好，你的"选择"是自主意志还是被引导的结果？AI Agent代你下单、回复消息时，责任与边界在哪里？
  \item \textbf{数据主权}：你的健康数据、脑电数据、生物特征数据，是属于你、属于设备厂商还是属于云服务商？《个人信息保护法》和欧盟GDPR只是起点，"数据可携带权"和"被遗忘权"的真正落地还有很长的路。
  \item \textbf{数字鸿沟}：AI手机、折叠屏、卫星通信都在拉大"用得起"和"用不起"之间的差距。数字基础设施的公平接入，将成为新的社会议题——就像20年前"会不会上网"一样关键。
  \item 当算法比你更懂你想看什么，你的"喜欢"有多少是出于自发？
  \item 当你每隔几分钟就不自觉地点亮屏幕，这是你的习惯还是被设计出来的"上瘾行为"？
\end{itemize}

\booksection{五、保持连接，也保持独立——你的数字生活哲学}

无论技术如何进化，下面这几条数字生活哲学，10年后依然适用：

\begin{enumerate}
  \item 做资源的分配者，而非被推送的接收者。
  \item 让自动化成为你忠实的"副驾驶"。
  \item 建立你稳固的"数字后方"。
  \item 定期执行"数字断舍离"与"社交斋戒"。
  \item 永远把安全和隐私作为一切便利的前提。
\end{enumerate}

\begin{tcbpractice}{来试试：制定你的《个人数字设备使用服务条款》}
为你和你的手机（以及未来的AR眼镜、AI助手）写一份简单的“服务条款”，来定义你们的关系。例如：

\begin{itemize}
  \item 第1条：我的注意力时间属于我。你在我工作/陪伴家人时应主动进入"勿扰"状态。
  \item 第2条：我的个人数据是我的财产。你的首要职责是安全加密与便捷同步。
  \item 第3条：我们的共同目标是提升生活效率与体验。若你通过诱导性设计让我沉迷，我有权将你卸载。
\end{itemize}

\end{tcbpractice}

\begin{tcbsummary}{本章小结}
\begin{itemize}
  \item Vision Pro与AR眼镜是手机之外的新计算入口，但短期难以取代手机
  \item 脑机接口距离消费级还有10-20年，伦理问题先于技术成熟
  \item 手机正向"隐身"方向进化，从设备中心走向身份中心
  \item AI伦理、数据主权、数字鸿沟是未来10年的核心议题
  \item 保持连接也保持独立，是数字生活不变的哲学
\end{itemize}

\end{tcbsummary}
